\begin{statement}[C.7.4; Proved conditionally: transcendence of the local norm]
Assume Grothendieck's period conjecture for $1$-motives over
$\overline\Q$. For every algebraic parameter $t\ne0,1$ at
which $P_t$ is nontorsion, $N(P_t)$ is transcendental.
More generally the conclusion holds for every nontorsion
algebraic point on an elliptic curve over $\overline\Q$.
The weaker named Relative Semi-abelian Conjecture 4.7 of
\citet{bertolin2022}, restricted to the groups and algebraic
lifts used in this proof, already suffices.
\end{statement}
\begin{proof}
Suppose $N(P)$ is algebraic. Statement C.7.2 supplies an
algebraic lift $R=(P,r)$ in $G_P(\C)^c$, and its conjugate
supplies a lift $\bar R$ in $G_{\bar P}(\C)^c$.
Let $v$ denote the logarithmic coordinate in
\eqref{eq:cp-v}. Put
\[
 K_0=\overline\Q\bigl(\text{entries of the full de Rham--Betti
                    period matrix of }G_P\times G_{\bar P}\bigr).
\]
Here ``full'' includes the integrals of differentials of the second
kind, hence the quasi-periods $\eta_j$, as well as the periods
$\omega_j$, the third-kind periods $\kappa_j$, and $2\pi i$.
It is the matrix $\Omega_G$ in their Definition 2.4, not merely
the lattice in the Lie algebra of the semi-abelian variety.
This field contains $u,\zeta(u)$ and their complex conjugates.
Indeed it contains the pure periods, $2\pi i$, and
$\kappa_j=\zeta(u)\omega_j-\eta_j u$; the determinant of
the two equations for $u,\zeta(u)$ is, up to sign,
$\omega_1\eta_2-\omega_2\eta_1= -2\pi i$.
This is also the observation of \citet[Remark 4.13]{bertolin2022}.
The real-linear expression for $\eta(u)$ is
\begin{equation}\label{eq:cp-real-linear}
 \eta(u)=
 \frac{(u\bar\omega_2-\bar u\omega_2)\eta_1
       +(\bar u\omega_1-u\bar\omega_1)\eta_2}
      {\omega_1\bar\omega_2-\bar\omega_1\omega_2}.
\end{equation}
The denominator is nonzero. Thus $u\eta(u)$ and its conjugate
belong to $K_0$. The compact equation $\chi(R)=0$ becomes
the algebraic relation over $K_0$
\begin{equation}\label{eq:cp-relation}
 v+\bar v
 =u\zeta(u)+\bar u\overline{\zeta(u)}
                -u\eta(u)-\bar u\overline{\eta(u)}.
\end{equation}
Complex conjugates here are coordinates of the conjugate
$1$-motive. They are not being treated as rational functions
of un-conjugated periods without justification.

Consider the period torsor of $\mathcal M$ over the period
torsor of $G=G_P\times G_{\bar P}$. The kernel of the
corresponding motivic groups is its weight $-2$ vector group
$\operatorname{Lie}Z$, of dimension $\dim Z$. There is no
weight $-1$ kernel: the point coordinates are exactly the
extension coordinates under the polarization, so $B_Q=0$.
More explicitly, projection of the smallest abelian variety $B$
to the variety generated by the extension coordinates is an
isogeny, since $B$ is the graph of the same point data in the
two factors. It therefore induces an isomorphism on the
weight $-1$ graded Lie algebras. The motive $[0\to G]$ has
only weights $-1$ and $-2$, and its unipotent Lie algebra has
no weight $-2$ graded piece. Thus the entire kernel is exactly
$W_{-2}\operatorname{Lie}G_{\rm mot}(\mathcal M)$; no additional
abelian direction or unaccounted logarithm remains in the fibre.
After evaluating the base periods, the fibre is an affine
torsor over $K_0$ under that vector group. Translation in
its two central coordinates changes $(v,\bar v)$ by
$2\pi i(c_1,c_2)$, with the signs fixed by using a positive
Tate loop for each motive. The linear functional
$c_1+c_2$ is nonzero on this central group. It is nonzero
on the whole plane when $\dim B=2$, and when $\dim B=1$
it is already nonzero on the bracket line
\eqref{eq:cp-positive-centre}, because $1+c>0$.

Consequently \eqref{eq:cp-relation} cuts a proper affine
hyperplane in the period-torsor fibre. It forces the relative
transcendence degree of the actual periods over $K_0$ to be
at most $\dim Z-1$. The pure-base period torsor has dimension
$\dim G_{\mathrm{mot}}(G)$, so the total transcendence degree
is at most $\dim G_{\mathrm{mot}}(\mathcal M)-1$.
This contradicts the asserted period conjecture. For real
data the single motive suffices: conjugation acts on its
Betti realization over $\Q$, and the same equation has a
nonzero coefficient on its one central direction. One can
also double the point and lift into the identity component,
choose $u,\omega_1$ real, and read the relation directly as
\[
 \omega_1 v-\omega_1u\zeta(u)+\eta_1u^2=0.
\]
Doubling preserves the equivalence of algebraicity of $N$
by the division formula and its algebraic nonzero factors.

For the weaker conjectural input, Conjecture 4.7 asserts
relative transcendence degree at least $2\dim B_Q+\dim Z$.
All $R_i$ used here are algebraic, $B_Q=0$, and the
generalized logarithms add exactly the central coordinates
to $K_0$. Thus it contradicts the same hyperplane relation
directly. Proposition 4.12 of that paper derives Conjecture
4.7 from the generalized period conjecture; in our case the
field of definition is algebraic, so the ordinary period
conjecture is enough. The weaker Conjecture 4.9, formulated
using arbitrary algebraic translates, gives a zero lower
bound at an algebraic point and does not by itself yield
this conclusion. Linear relations among $1$-periods are
also insufficient: after clearing the denominator in
\eqref{eq:cp-real-linear}, the prohibited relation is
polynomial, not a linear relation with algebraic coefficients.
\end{proof}
